\documentclass[journal]{IEEEtran}

% --- packages ---
\usepackage{amsmath,amssymb}
\usepackage{graphicx}
\usepackage{booktabs}
\usepackage{url}
\usepackage{textcomp}
\usepackage[hidelinks]{hyperref}

% Figures live in ../results ; adjust if you copy them next to this file.
\graphicspath{{../results/}{./}}

\begin{document}

\title{A Parsimonious, Externally Validated Additive Survival Model for Lung
Adenocarcinoma: When Simplicity Generalizes Better Than an Ensemble}

\author{
  \IEEEauthorblockN{Author Name}%
  \thanks{Manuscript prepared \today. Author is with the Department of
  \emph{[affiliation]}, \emph{[institution]} (e-mail: \emph{[email]}).}%
}

\markboth{IEEE Journal, Draft}{}

\maketitle

\begin{abstract}
Prognostic models for lung adenocarcinoma (LUAD) are frequently developed on a
single cohort and are seldom validated externally or assessed for
interpretability, which limits their clinical usefulness. We ask a narrow but
practical question: does added model complexity translate into better
\emph{generalizable} survival prediction? Using the TCGA-LUAD cohort
($n=501$; $181$ deaths), we benchmarked four survival learners---random
survival forest (RSF), gradient boosting survival (GBS), XGBoost with a Cox
objective, and DeepSurv---together with a spline generalized-additive-model
(GAM) stacking ensemble, against a single, transparent feature-level GAM
(elastic-net Cox with cubic B-spline smooths on age, tumor mutational burden,
mutation count, and fraction of genome altered; stage and sex entered
linearly). Every model was restricted to the six clinical and genomic features
shared with an independent external cohort (OncoSG, $n=302$; $94$ deaths), so
that the identical model could be transported. We evaluated discrimination with
leakage-free nested cross-validation and bootstrap confidence intervals,
performed Kaplan--Meier risk stratification, and tested for nonlinearity with a
likelihood-ratio (LR) test. Internally the ensemble was competitive, but on
external validation it degraded sharply (concordance index, $C=0.595$),
whereas the feature-level GAM transferred without loss ($C=0.660$ internal;
$C=0.679$ external) and matched the best individual learners (XGBoost $0.675$,
DeepSurv $0.674$). The GAM separated external risk tertiles with a log-rank
$p=1.4\times10^{-6}$. The LR test confirmed that nonlinearity improved fit
($\text{LR}=17.6$, $\mathrm{df}=8$, $p=0.024$), driven chiefly by a J-shaped age
effect ($p=0.032$). A parsimonious, interpretable additive model generalized
more reliably than a complex ensemble that overfit the development cohort,
arguing for interpretable simplicity in cross-cohort LUAD prognosis using
routinely available features.
\end{abstract}

\begin{IEEEkeywords}
Survival analysis, lung adenocarcinoma, generalized additive model, external
validation, model interpretability, concordance index, TCGA.
\end{IEEEkeywords}

% =====================================================================
\section{Introduction}
% =====================================================================
\IEEEPARstart{L}{ung} adenocarcinoma (LUAD) is the most common histological
subtype of lung cancer and a leading cause of cancer death worldwide~\cite{siegel2023}.
Accurate estimates of a patient's prognosis inform surveillance and treatment
decisions, and a large body of work has applied machine learning to survival
prediction from clinical and molecular data. Much of that work, however, shares
two recurring weaknesses. First, models are typically reported on a single
cohort, using internal cross-validation or a held-out split, without
validation on a genuinely independent population; internal performance is a
weak proxy for how a model will behave elsewhere~\cite{tripod2015}. Second, the
better-performing methods tend to be opaque---tree ensembles, gradient boosting,
deep networks, or stacked combinations of them---so that even when they
discriminate well, they offer little insight a clinician can inspect or trust.

We were motivated by a simple observation while building such a system. A
carefully engineered stacking ensemble that combined four survival learners
through a GAM meta-model performed respectably on the TCGA-LUAD cohort, yet we
had no evidence it would hold up on new patients. When we obtained an
independent cohort and tested it, the ensemble's advantage evaporated. That
failure prompted the central question of this paper: for cross-cohort LUAD
prognosis, is the complexity actually buying us anything?

To answer it fairly we did two things. We restricted every model to the set of
features that both our development and validation cohorts actually record, so
that the same model---not a re-fitted approximation---could be transported. And
we contrasted the ensemble and its constituent learners against a deliberately
simple alternative: a feature-level generalized additive model that keeps the
familiar Cox proportional-hazards form but allows smooth, nonlinear effects for
continuous covariates~\cite{cox1972,hastie1990}. The additive model has the
practical virtues the ensemble lacks---each covariate's contribution can be
plotted and read directly---so if it also generalizes at least as well, the
case for the ensemble becomes hard to make.

Our contributions are: (i) a leakage-free benchmarking protocol that compares
five modeling strategies on an identical, transportable feature set with
external validation on an independent cohort; (ii) evidence that a stacking
ensemble which looks competitive internally can overfit and fail to
generalize, while a parsimonious additive model transfers without loss; and
(iii) a formal likelihood-ratio test that isolates \emph{where} nonlinearity is
statistically justified, so that interpretability claims rest on a test rather
than on the appearance of a curve.

% =====================================================================
\section{Materials and Methods}
% =====================================================================

\subsection{Cohorts}
The development cohort was the TCGA lung adenocarcinoma project
(\texttt{luad\_tcga\_pan\_can\_atlas\_2018})~\cite{tcga2014}, obtained through
cBioPortal~\cite{cbioportal2012}. After removing patients with missing or
non-positive overall-survival time, $501$ patients remained, of whom $181$
($36.1\%$) had a recorded death. The external validation cohort was the OncoSG
East-Asian lung adenocarcinoma study (\texttt{luad\_oncosg\_2020})~\cite{chen2020},
also from cBioPortal, comprising $302$ evaluable patients with $94$ deaths
($31.1\%$). The two cohorts differ in ancestry, geography, and sequencing
platform, which makes the external test a demanding one.

\subsection{Features and harmonization}
Only features present in \emph{both} cohorts were used, so that the fitted model
could be applied unchanged to the external data. Six features met this
criterion: age at diagnosis, sex, AJCC tumor stage, tumor mutational burden
(TMB, nonsynonymous), somatic mutation count, and fraction of genome altered
(FGA). Because the two cohorts encode stage at different granularity, we
collapsed both to the four major stages (I--IV). Sex and survival-status
encodings were reconciled between cohorts. We emphasize that this restriction is
a requirement of honest external validation rather than a performance-driven
feature selection; features available only in TCGA (e.g., hypoxia scores,
aneuploidy score, microsatellite-instability metrics) were necessarily
excluded.

\subsection{Preprocessing}
All preprocessing was fit on TCGA alone and applied to OncoSG, to avoid
information leakage from the test cohort. Categorical variables (stage, sex)
were one-hot encoded with the first level dropped; missing continuous values
were median-imputed; and features were standardized. TMB and mutation count,
which are strongly right-skewed, were $\log(1+x)$-transformed before
standardization. For the additive model, external spline inputs were clipped to
the training knot range so that the basis was never evaluated outside its
support.

\subsection{Models}
We compared six strategies. The four base learners were: a random survival
forest~\cite{ishwaran2008} ($500$ trees, $\sqrt{p}$ features per split, minimum
leaf size $3$); gradient boosting survival analysis with regression-tree base
learners ($500$ stages, learning rate $0.05$, tree depth $3$); XGBoost~\cite{chen2016} with the
\texttt{survival:cox} objective and early stopping; and DeepSurv~\cite{katzman2018},
a Cox proportional-hazards neural network ($128$--$64$--$1$ units, ReLU
activations, dropout $0.3$, Adam optimization with early stopping). For
survival:cox we encoded censoring through the sign of the label so that
censored patients contribute to the risk sets rather than being discarded.

The fifth strategy was a \emph{stacking ensemble}: out-of-fold predictions from
the four base learners were combined by a GAM meta-learner (an elastic-net Cox
model with cubic B-spline smooths on each base-learner score). To prevent
optimistic bias, base-learner early stopping within each fold used a nested
internal split, never the held-out fold on which predictions were formed.

The sixth strategy---our proposed model---was a \emph{feature-level additive
survival model}. It is an elastic-net Cox model ($\ell_1$ ratio $0.5$) in which
each continuous feature enters through a cubic B-spline basis (four degrees of
freedom), and stage and sex enter linearly:
\begin{equation}
\log h_i(t) = \log h_0(t) + \sum_{k} f_k(x_{ik}) + \boldsymbol{\beta}^{\top}\mathbf{z}_i,
\end{equation}
where $f_k$ are the estimated smooths for age, TMB, mutation count, and FGA, and
$\mathbf{z}_i$ collects the categorical terms. The regularization strength was
chosen by internal cross-validation. This is the only model whose effect of
each covariate can be read directly from a plot.

\subsection{Evaluation}
Discrimination was measured by Harrell's concordance index
($C$)~\cite{harrell1996}. Internal performance on TCGA used nested five-fold
cross-validation for the additive model and five-fold out-of-fold predictions
for the base learners. External performance used the model fit on all of TCGA
and applied once to OncoSG, with $95\%$ confidence intervals from $1{,}000$
bootstrap resamples of the external cohort. We stratified patients into risk
tertiles and compared survival with the log-rank test and Kaplan--Meier curves.
Finally, to test whether the additive model's flexibility was warranted, we
performed a likelihood-ratio test comparing the full spline model against a
model in which the same continuous covariates entered linearly, using the Cox
partial likelihood; we report an overall test and a per-feature test. All
experiments used a fixed random seed. Analyses were implemented in Python with
scikit-survival~\cite{polsterl2020}.

% =====================================================================
\section{Results}
% =====================================================================

\subsection{External validation reverses the internal ranking}
Table~\ref{tab:cindex} reports discrimination on both cohorts. Internally, the
base learners were unremarkable ($C$ between $0.605$ and $0.650$) and the
full-feature stacking ensemble reached a held-out $C=0.678$. The decisive
comparison, however, is the external column. The feature-level additive model
achieved $C=0.679$ ($95\%$ CI $0.623$--$0.735$) on OncoSG---essentially
identical to its internal estimate of $0.660$, indicating no loss under cohort
shift. The individual learners transported comparably (XGBoost $0.675$,
DeepSurv $0.674$, RSF $0.660$). The stacking ensemble was the sole exception: it
fell to $C=0.595$ ($0.530$--$0.660$), well below every model it was built from.
In other words, the extra machinery of stacking did not merely fail to help---it
actively hurt generalization, because the meta-learner had fit relationships
specific to the development cohort.

\begin{table}[t]
\caption{Discrimination (concordance index) on TCGA (internal,
cross-validated) and OncoSG (external), using the six common features. External
values include $95\%$ bootstrap confidence intervals.}
\label{tab:cindex}
\centering
\begin{tabular}{lcc}
\toprule
Model & TCGA (internal) & OncoSG (external) \\
\midrule
Feature-level GAM (proposed) & $0.660$ & $\mathbf{0.679}$ $[0.623,0.735]$ \\
XGBoost--Cox                 & $0.605$ & $0.675$ $[0.620,0.731]$ \\
DeepSurv                     & $0.650$ & $0.674$ $[0.614,0.733]$ \\
Random survival forest       & $0.621$ & $0.660$ $[0.600,0.721]$ \\
Gradient boosting survival   & $0.635$ & $0.641$ $[0.577,0.702]$ \\
GAM stacking ensemble        & --\,\textsuperscript{a} & $0.595$ $[0.530,0.660]$ \\
\bottomrule
\end{tabular}
\\[2pt]
\raggedright
\footnotesize\textsuperscript{a}The stacking ensemble was assessed internally
with the full $25$-feature TCGA set (held-out $C=0.678$, $[0.561,0.775]$); on
the six common features it is shown for external comparison only.
\end{table}

\subsection{Risk stratification holds in the external cohort}
Stratifying OncoSG patients into tertiles of the additive model's risk score
produced well-separated survival curves (Fig.~\ref{fig:km}), with a log-rank
$p=1.4\times10^{-6}$ between the low- and high-risk groups (multivariable
log-rank $p=4.6\times10^{-6}$). The same stratification on TCGA gave
$p=6.4\times10^{-10}$. For contrast, tertiles from the stacking ensemble
separated the external cohort far more weakly ($p=2.3\times10^{-3}$), consistent
with its diminished discrimination.

\begin{figure}[t]
\centering
\includegraphics[width=\columnwidth]{gam14_km_oncosg.png}
\caption{Kaplan--Meier curves for OncoSG patients stratified into tertiles of
the feature-level additive model's risk score. Groups are well separated on this
fully external cohort (log-rank $p=1.4\times10^{-6}$).}
\label{fig:km}
\end{figure}

\subsection{Where nonlinearity is (and is not) justified}
The additive model's smooths are shown in Fig.~\ref{fig:smooths}. The
likelihood-ratio test (Table~\ref{tab:lr}) confirms that allowing nonlinearity
improves fit over a linear Cox model overall ($\text{LR}=17.6$, $\mathrm{df}=8$,
$p=0.024$). Decomposed by feature, the effect is carried mainly by age
($p=0.032$), whose contribution is J-shaped---risk is lowest in middle age and
rises in the elderly---a pattern a linear term cannot represent. Mutation count
was borderline ($p=0.054$). We note candidly that the visually striking curves
for TMB and FGA are \emph{not} statistically supported ($p=0.18$ and $p=0.26$);
for these features a linear term is adequate, and we caution against
over-interpreting their smooths. This is precisely the kind of distinction the
test is meant to enforce.

\begin{figure}[t]
\centering
\includegraphics[width=\columnwidth]{gam14_smooths.png}
\caption{Estimated smooth contributions $f_k$ (solid) with a linear reference
(dashed) for each continuous feature, fit on TCGA. Departure from the dashed
line indicates nonlinearity. Only the age effect is statistically significant
(Table~\ref{tab:lr}); rug marks show data density, and the sparse tails should
be read with caution.}
\label{fig:smooths}
\end{figure}

\begin{table}[t]
\caption{Likelihood-ratio test of nonlinearity (spline vs.\ linear) on TCGA,
using the Cox partial likelihood.}
\label{tab:lr}
\centering
\begin{tabular}{lccc}
\toprule
Term & LR statistic & df & $p$ \\
\midrule
Overall (spline vs.\ linear) & $17.63$ & $8$ & $\mathbf{0.024}$ \\
Age                          & $6.89$  & $2$ & $\mathbf{0.032}$ \\
Mutation count               & $5.82$  & $2$ & $0.054$ \\
TMB                          & $3.48$  & $2$ & $0.176$ \\
Fraction genome altered      & $2.69$  & $2$ & $0.261$ \\
\bottomrule
\end{tabular}
\end{table}

% =====================================================================
\section{Discussion}
% =====================================================================
The pattern across Table~\ref{tab:cindex} is the paper's main message. Every
model except the stacking ensemble carried its performance from TCGA to OncoSG;
several even discriminated slightly better on the external cohort, which argues
against any of them having overfit. The ensemble alone collapsed. Its
meta-learner had been trained on the joint distribution of the base-learner
scores in TCGA, and that geometry does not survive a change of cohort and
sequencing platform---so the very step intended to squeeze out extra
performance became the step that failed to transport. This is a concrete,
cautionary instance of a familiar risk in stacked
generalization~\cite{wolpert1992}: added flexibility at the combination stage
can encode cohort-specific idiosyncrasies.

Against that backdrop, the feature-level additive model is attractive not
because it wins by a wide margin---on external data it is within the confidence
interval of XGBoost and DeepSurv---but because it matches them while being
transparent and stable. Its risk score is a sum of inspectable functions; its
nonlinearity is limited to where a formal test supports it; and it uses only
features a clinic would already have. For a prognostic tool intended to travel
between institutions, those properties are worth more than a fractional,
non-significant gain in $C$-index from an opaque model.

Our study has clear limitations. It uses a single external cohort, and although
that cohort is genuinely independent, a second would strengthen the
generalizability claim. The model is deliberately narrow---six routinely
available features---and richer molecular data (gene expression, driver-gene
status, or the TCGA-only genomic scores we excluded) may well add prognostic
signal; we traded that potential signal for transportability. The event count
($181$ deaths) limits statistical power, which is visible in the wide confidence
intervals and in the borderline nonlinearity of mutation count. Finally, TMB and
mutation count are measured on different platforms across cohorts, introducing
distributional shift in exactly the features whose nonlinearity we could not
confirm; that the model transported regardless is reassuring, but the point
deserves a reader's attention.

% =====================================================================
\section{Conclusion}
% =====================================================================
On a demanding external test, a parsimonious additive survival model for LUAD
generalized as well as the strongest individual machine-learning learners and
markedly better than a stacking ensemble that had looked competitive
internally. The additive model is interpretable by construction, its
nonlinearity is confined to where a likelihood-ratio test justifies it, and it
depends only on features available across cohorts. For cross-cohort prognosis,
these results favor interpretable simplicity over ensemble complexity, and they
underline that external validation---not internal performance---is what should
decide between competing models.

% =====================================================================
\section*{Reproducibility}
% =====================================================================
All preprocessing was fit on the development cohort only; a fixed random seed
was used throughout. Analysis code is available at
\url{https://github.com/userlk327/tcga-luda-gam}. Both cohorts are publicly
available through cBioPortal.

% =====================================================================
\begin{thebibliography}{99}
% =====================================================================
\bibitem{siegel2023} R.~L. Siegel, K.~D. Miller, N.~S. Wagle, and A.~Jemal,
``Cancer statistics, 2023,'' \emph{CA Cancer J. Clin.}, vol.~73, no.~1,
pp.~17--48, 2023.

\bibitem{tripod2015} G.~S. Collins, J.~B. Reitsma, D.~G. Altman, and K.~G.~M.
Moons, ``Transparent reporting of a multivariable prediction model for
individual prognosis or diagnosis (TRIPOD),'' \emph{Ann. Intern. Med.},
vol.~162, no.~1, pp.~55--63, 2015.

\bibitem{cox1972} D.~R. Cox, ``Regression models and life-tables,'' \emph{J. R.
Stat. Soc. B}, vol.~34, no.~2, pp.~187--220, 1972.

\bibitem{hastie1990} T.~Hastie and R.~Tibshirani, \emph{Generalized Additive
Models}. London, U.K.: Chapman \& Hall, 1990.

\bibitem{tcga2014} Cancer Genome Atlas Research Network, ``Comprehensive
molecular profiling of lung adenocarcinoma,'' \emph{Nature}, vol.~511,
no.~7511, pp.~543--550, 2014.

\bibitem{cbioportal2012} E.~Cerami \emph{et~al.}, ``The cBio cancer genomics
portal: An open platform for exploring multidimensional cancer genomics data,''
\emph{Cancer Discov.}, vol.~2, no.~5, pp.~401--404, 2012.

\bibitem{chen2020} J.~Chen \emph{et~al.}, ``Genomic landscape of lung
adenocarcinoma in East Asians,'' \emph{Nat. Genet.}, vol.~52, no.~2,
pp.~177--186, 2020.

\bibitem{ishwaran2008} H.~Ishwaran, U.~B. Kogalur, E.~H. Blackstone, and M.~S.
Lauer, ``Random survival forests,'' \emph{Ann. Appl. Stat.}, vol.~2, no.~3,
pp.~841--860, 2008.

\bibitem{chen2016} T.~Chen and C.~Guestrin, ``XGBoost: A scalable tree boosting
system,'' in \emph{Proc. 22nd ACM SIGKDD Int. Conf. Knowl. Discovery Data
Mining}, 2016, pp.~785--794.

\bibitem{katzman2018} J.~L. Katzman \emph{et~al.}, ``DeepSurv: Personalized
treatment recommender system using a Cox proportional hazards deep neural
network,'' \emph{BMC Med. Res. Methodol.}, vol.~18, no.~1, p.~24, 2018.

\bibitem{polsterl2020} S.~P\"olsterl, ``scikit-survival: A library for
time-to-event analysis built on top of scikit-learn,'' \emph{J. Mach. Learn.
Res.}, vol.~21, no.~212, pp.~1--6, 2020.

\bibitem{harrell1996} F.~E. Harrell, K.~L. Lee, and D.~B. Mark, ``Multivariable
prognostic models: Issues in developing models, evaluating assumptions and
adequacy, and measuring and reducing errors,'' \emph{Stat. Med.}, vol.~15,
no.~4, pp.~361--387, 1996.

\bibitem{wolpert1992} D.~H. Wolpert, ``Stacked generalization,'' \emph{Neural
Netw.}, vol.~5, no.~2, pp.~241--259, 1992.
\end{thebibliography}

\end{document}
